% !TeX program = pdflatex
\documentclass[11pt,a4paper]{article}

% Compile with pdfLaTeX:
%   latexmk -pdf legged_mppi_analysis_draft_en.tex
\usepackage[utf8]{inputenc}
\usepackage[T1]{fontenc}
\usepackage{amsmath,amssymb,amsthm}
\usepackage{bm}
\usepackage{geometry}
\usepackage{booktabs}
\usepackage{array}
\usepackage{algorithm}
\usepackage{algorithmic}
\usepackage{xcolor}
\usepackage{hyperref}

\geometry{margin=2.4cm}
\hypersetup{
  colorlinks=true,
  linkcolor=blue!55!black,
  urlcolor=blue!55!black,
  citecolor=blue!55!black
}

\newcommand{\R}{\mathbb{R}}
\newcommand{\norm}[1]{\left\lVert #1 \right\rVert}
\newcommand{\wnorm}[2]{\left\lVert #1 \right\rVert_{#2}^2}
\newcommand{\clip}{\operatorname{clip}}
\newcommand{\diag}{\operatorname{diag}}
\newcommand{\code}[1]{\texttt{#1}}

\title{\code{legged\_mppi}: Whole-Body Locomotion MPPI\\
Mathematical Formulation and Implementation Analysis}
\author{Notes based on analysis of the \code{legged\_mppi} codebase}
\date{\today}

\begin{document}
\maketitle
\tableofcontents
\newpage

%=============================================================================
\section{Scope and Executive Summary}
%=============================================================================

This document analyzes the locomotion controller implemented by
\code{legged\_mppi}.  The presentation follows the problem-oriented style of
the accompanying OCS2 legged-MPC note, but the mathematical description here
is derived from the actual Python, YAML, gait, and MuJoCo files in this
repository.  The primary implementation sources are:
\begin{itemize}
  \item \code{control/controllers/base\_controller.py},
  \item \code{control/controllers/mppi\_locomotion.py},
  \item \code{control/gait\_scheduler/scheduler.py},
  \item \code{utils/tasks.py} and the \code{mppi\_gait\_config\_*.yml} files,
  \item the Go1 MJCF files under \code{models/go1}, and
  \item \code{interface/simulator.py} and
        \code{scripts/run\_mppi\_locomotion.py}.
\end{itemize}

The main architectural choices are:
\begin{enumerate}
  \item The controller uses the complete 37-dimensional floating-base Go1
        state rather than a reduced centroidal state.
  \item The optimized input is a sequence of 12 desired joint positions, not
        joint torques, joint velocities, or contact forces.
  \item Nonlinear rigid-body dynamics and contact are evaluated by parallel
        MuJoCo rollouts.  No analytical dynamics derivatives are required.
  \item A periodic, precomputed 24-row gait file supplies desired joint
        positions and velocities over the prediction horizon.
  \item Each rollout is scored by body-state tracking, gait tracking, and a
        virtual PD-effort penalty.
  \item The implemented update is an MPPI-style exponentially weighted mean
        of sampled absolute action sequences.  It is not the canonical
        path-integral noise-increment update.
\end{enumerate}

Contact is therefore \emph{implicit in the rollout model}.  The gait file is
a kinematic motion prior; it does not impose a hard stance/swing contact-mode
schedule.  Feasibility is produced by action clipping, MuJoCo joint and force
limits, and the simulated contact dynamics rather than explicit OCP
constraints.

%=============================================================================
\section{Robot Model, State, and Input}
%=============================================================================

\subsection{Degrees of Freedom and Joint Order}

The locomotion model is a floating-base Unitree Go1 with twelve actuated
joints.  Its MuJoCo dimensions are
\begin{equation}
  n_q=19, \qquad n_v=18, \qquad n_u=12.
\end{equation}
The actuator order used consistently by the controller is
\begin{align}
\bm{q}_j = [&q_{\mathrm{FR,h}},q_{\mathrm{FR,t}},q_{\mathrm{FR,c}},
q_{\mathrm{FL,h}},q_{\mathrm{FL,t}},q_{\mathrm{FL,c}}, \nonumber\\
&q_{\mathrm{RR,h}},q_{\mathrm{RR,t}},q_{\mathrm{RR,c}},
q_{\mathrm{RL,h}},q_{\mathrm{RL,t}},q_{\mathrm{RL,c}}]^\top ,
\end{align}
where $\mathrm{h}$, $\mathrm{t}$, and $\mathrm{c}$ denote hip abduction,
thigh, and calf joints, respectively.

\subsection{Observation State}

The controller observation concatenates MuJoCo \code{qpos} and \code{qvel}:
\begin{equation}
  \boxed{
  \bm{x}=
  \begin{bmatrix}
    \bm{p}_b\\
    \bm{q}_b\\
    \bm{q}_j\\
    \bm{v}_b\\
    \bm{\omega}_b\\
    \dot{\bm{q}}_j
  \end{bmatrix}
  \in\R^{37}.}
  \label{eq:state}
\end{equation}
Here $\bm{p}_b\in\R^3$ is the floating-base position in the world frame,
$\bm{q}_b=[q_w,q_x,q_y,q_z]^\top\in\mathbb{S}^3$ is its quaternion,
$\bm{q}_j\in\R^{12}$ is the joint position, $\bm{v}_b$ and
$\bm{\omega}_b$ are base linear and angular velocities, and
$\dot{\bm{q}}_j\in\R^{12}$ contains joint velocities.

\begin{center}
\begin{tabular}{ccl}
\toprule
Python slice & Dimension & Quantity\\
\midrule
$0{:}3$   & 3  & base position $\bm{p}_b$\\
$3{:}7$   & 4  & base quaternion $\bm{q}_b$\\
$7{:}19$  & 12 & joint position $\bm{q}_j$\\
$19{:}22$ & 3  & base linear velocity $\bm{v}_b$\\
$22{:}25$ & 3  & base angular velocity $\bm{\omega}_b$\\
$25{:}37$ & 12 & joint velocity $\dot{\bm{q}}_j$\\
\bottomrule
\end{tabular}
\end{center}

MuJoCo's \code{FULLPHYSICS} rollout state additionally stores time:
\begin{equation}
  \bm{x}^{\mathrm{full}}=
  [t,\bm{q}_{\mathrm{pos}}^\top,\bm{q}_{\mathrm{vel}}^\top]^\top
  \in\R^{38}.
\end{equation}
The implementation removes the first element before evaluating cost.

\subsection{Control Input and Low-Level Actuation}

The optimized input is a desired joint-position vector:
\begin{equation}
  \boxed{\bm{u}=\bm{q}_{j,\mathrm{des}}\in\R^{12}},
  \qquad \bm{u}_{\min}\leq\bm{u}\leq\bm{u}_{\max}.
  \label{eq:input}
\end{equation}
The limits are repeated once per leg:
\begin{equation}
  \bm{u}_{\min}^{\mathrm{leg}}=
  [-0.863,-0.686,-2.818]^\top,
  \quad
  \bm{u}_{\max}^{\mathrm{leg}}=
  [0.863,4.501,-0.888]^\top.
\end{equation}

In the default rollout MJCF, an affine position actuator approximately applies
\begin{equation}
  \tau_i=operatorname{sat}_{[\tau_i^-,\tau_i^+]}
  \left(55(u_i-q_{j,i})-3\dot q_{j,i}\right).
  \label{eq:mjcf_pd}
\end{equation}
The nominal force range is $[-23.7,23.7]$ N m, enlarged to
$[-35.55,35.55]$ N m for the calf.  The hardware interface sends the same
desired positions in Unitree position-control mode with $K_p=55$, $K_d=3$,
zero desired velocity, and zero feedforward torque.

%=============================================================================
\section{Prediction Dynamics and Contact}
%=============================================================================

\subsection{Discrete MuJoCo Dynamics}

For analysis, one rollout step can be written as
\begin{equation}
  \bm{x}_{k+1}=F_{\mathrm{MJ}}
  (\bm{x}_k,\bm{u}_k;\mathcal{M},\Delta t,\bm{\eta}),
  \label{eq:mujoco_dynamics}
\end{equation}
where $\mathcal{M}$ is the MJCF robot/environment model and $\bm{\eta}$
collects numerical and contact parameters.  The default values include
\begin{equation}
  \Delta t=0.01\ \mathrm{s},\qquad
  \code{o\_solref}=[0.02,1.0],
\end{equation}
with pyramidal friction cones in the prediction model.

Equation~\eqref{eq:mujoco_dynamics} includes floating-base multibody
dynamics, actuator saturation, joint limits, gravity, collision detection,
normal contact, and friction.  Unlike the SRBD model in the OCS2 note, the
implementation does not expose contact force as an optimization variable and
does not explicitly write a centroidal momentum equation.

For each controller update, all $N$ candidates begin from the repeated current
observation $\hat{\bm{x}}_t$ and are integrated for $H$ steps.  The output
tensor has shape
\begin{equation}
  N\times H\times38,
\end{equation}
and is evaluated as $N\times H\times37$ after removing time.

\subsection{Parallel Rollout Execution}

Candidates are divided into chunks and submitted to a thread pool.  Every
worker owns a thread-local \code{MjData}, while all workers share the immutable
\code{MjModel}.  The MuJoCo rollout API performs the $H$-step integrations.
For the default configuration this means
\begin{equation}
  N H = 30\times40=1200
\end{equation}
contact-dynamics steps per update, distributed over five workers.

\subsection{Contact Treatment}

The controller does not optimize a binary contact mode
$\sigma_k$, and the gait file does not activate mode-dependent constraints.
Instead,
\begin{equation}
  \text{desired joint gait}
  \longrightarrow \text{foot motion}
  \longrightarrow \text{MuJoCo collision/contact response}.
\end{equation}
Thus, stance and swing emerge from tracking a periodic joint-space reference
inside the simulated contact dynamics.  Contact timing can deviate from the
nominal gait when sampled actions, terrain, or tracking errors change the
motion.

The prediction and closed-loop simulator normally load separate MJCF files.
For the default walking task these are \code{go1\_scene\_mppi\_pyr.xml} for
rollouts and \code{go1\_scene\_mppi.xml} for execution.  They share the same
Go1 structure and actuator convention, but the prediction variant explicitly
selects pyramidal contact.  Any further mismatch in geometry, contacts, or
parameters acts as model error.

%=============================================================================
\section{Finite-Horizon Stochastic Optimization Problem}
%=============================================================================

At update time $t$, let
\begin{equation}
  \bm{U}=[\bm{u}_0,\ldots,\bm{u}_{H-1}]\in\R^{H\times12}
\end{equation}
be the warm-start action sequence.  The controller samples $N$ bounded
candidate sequences $\bm{U}^{(n)}$ and approximately solves
\begin{equation}
  \boxed{
  \begin{aligned}
  \min_{\bm{U}}\quad &
    S(\bm{U};\hat{\bm{x}}_t,\bm{X}^{\mathrm{ref}}_t)
    =\sum_{k=0}^{H-1}
      \ell(\bm{x}_{k+1},\bm{u}_k,\bm{x}^{\mathrm{ref}}_k),\\
  \text{s.t.}\quad &
    \bm{x}_{k+1}=F_{\mathrm{MJ}}(\bm{x}_k,\bm{u}_k),\\
  & \bm{x}_0=\hat{\bm{x}}_t,\\
  & \bm{u}_{\min}\leq\bm{u}_k\leq\bm{u}_{\max}.
  \end{aligned}}
  \label{eq:stochastic_ocp}
\end{equation}
There is no separately weighted terminal cost in the locomotion
implementation.  The rollout sum includes all returned states with the same
stage-cost expression.

This is a sample-based numerical optimization, not a direct transcription of
the dynamics.  States are never independent decision variables, and dynamics
or cost Jacobians are not computed.

%=============================================================================
\section{Reference Generation and Gait Scheduling}
%=============================================================================

\subsection{Waypoint-Level Body Reference}

For waypoint $g$, the task dictionary defines a desired position, nominal
orientation, commanded planar velocity, distance threshold, gait name, and
waiting time.  The controller stores
\begin{equation}
  \bm{r}_b^{\mathrm{ref}}=
  \begin{bmatrix}
    \bm{p}_g\\
    \bm{q}_g\\
    v_{x,g}^{\mathrm{cmd}}\\
    v_{y,g}^{\mathrm{cmd}}\\
    0\\0\\0\\0
  \end{bmatrix}\in\R^{13}.
\end{equation}
The final four zeros correspond to desired vertical linear velocity and three
angular-velocity components.

When the position error exceeds $0.1$ m and the waypoint timer is not in its
waiting state, $\bm{q}_g$ is recomputed to point toward the goal.  If
$\bm{d}=\bm{p}_g-\bm{p}_b$, then
\begin{align}
  \psi_g &= \operatorname{atan2}(d_y,d_x),\\
  \theta_g &= -\operatorname{atan2}
  \left(d_z,\sqrt{d_x^2+d_y^2}\right),\\
  \bm{q}_g &= \bm{q}_z(\psi_g)\otimes\bm{q}_y(\theta_g).
\end{align}
Otherwise, the desired orientation is the identity quaternion.

\subsection{Periodic Joint-Space Gait}

Each gait is loaded from a tab-separated matrix
\begin{equation}
  \bm{G}\in\R^{24\times L}.
\end{equation}
The first twelve rows are joint-position references and the last twelve rows
are joint-velocity references:
\begin{equation}
  \bm{G}=
  \begin{bmatrix}
    \bm{G}_q\\
    \bm{G}_{\dot q}
  \end{bmatrix}.
\end{equation}
The \code{in\_place} and \code{walk\_fast} files have $L=100$ columns,
whereas the two medium-speed files used by \code{walk} and \code{trot} have
$L=74$.  At controller phase $s$, the reference horizon is selected cyclically:
\begin{equation}
  \bm{G}_{s,H}=
  [\bm{G}_{:,s},\bm{G}_{:,s+1},\ldots,
  \bm{G}_{:,s+H-1}],
\end{equation}
where column indices wrap modulo $L$.  After every MPPI update, the index
array is rolled by one column.

The active scheduler is selected from
\begin{equation}
  \{\code{in\_place},\code{walk},\code{walk\_fast},\code{trot}\}
\end{equation}
when the waypoint changes.  Each scheduler is a persistent object, so
reselecting a gait resumes its previously stored phase rather than
automatically resetting phase zero.

\subsection{Complete State Reference}

For rollout step $k$, the reference used by the cost is
\begin{equation}
  \boxed{
  \bm{x}^{\mathrm{ref}}_k=
  \begin{bmatrix}
    \bm{p}_g\\
    \bm{q}_g\\
    \bm{G}_{q,s+k}\\
    v_{x,g}^{\mathrm{cmd}}\\
    v_{y,g}^{\mathrm{cmd}}\\
    0\\
    \bm{0}_3\\
    \bm{G}_{\dot q,s+k}
  \end{bmatrix}\in\R^{37}.}
  \label{eq:state_reference}
\end{equation}
Before cost evaluation, the predicted world-frame base linear velocity is
rotated into the predicted body frame:
\begin{equation}
  {}^B\bm{v}_{b,k}=R(\bm{q}_{b,k})^\top{}^W\bm{v}_{b,k}.
\end{equation}
Accordingly, the commanded $x$ and $y$ velocities in
Eq.~\eqref{eq:state_reference} are interpreted in the body frame.  Base
angular velocity is not transformed by this cost-side operation.

%=============================================================================
\section{Sampling and MPPI-Style Update}
%=============================================================================

\subsection{Warm Start}

At initialization the entire horizon is filled with
\begin{equation}
\bm{u}_{\mathrm{init}}=
\bm{1}_2\otimes
[-0.3,1.34,-2.83,\ 0.3,1.34,-2.83]^\top,
\end{equation}
which produces the twelve entries for the four legs.  After an update, the
new plan is shifted:
\begin{equation}
  \bm{u}_{k,t+1}^{\mathrm{nom}}=
  \begin{cases}
    \bm{u}_{k+1,t}^{+}, & k=0,\ldots,H-2,\\
    \bm{u}_{H-1,t}^{+}, & k=H-1.
  \end{cases}
  \label{eq:warm_start}
\end{equation}
The last optimized action is therefore held at the tail.

\subsection{Cubic Gaussian Sampling}

The default sampler uses $K=4$ knot indices.  For $H=40$, rounded linearly
spaced indices give
\begin{equation}
  \mathcal{K}=\{0,13,26,39\}.
\end{equation}
For sample $n$, independent Gaussian knot noise is generated as
\begin{equation}
  \bm{\epsilon}_{n,j}\sim
  \mathcal{N}(\bm{0},\diag(\bm{\sigma}^2)),
  \qquad j\in\mathcal{K}.
\end{equation}
The default walking standard deviation repeats
\begin{equation}
  \bm{\sigma}^{\mathrm{leg}}=[0.06,0.1,0.1]^\top
\end{equation}
for all four legs.  During a \code{trot} task phase, the controller changes
this to $[0.06,0.2,0.2]^\top$ per leg.  A cubic spline interpolates the noisy
knot actions across the full horizon:
\begin{equation}
  \tilde{\bm{u}}^{(n)}_k=
  \operatorname{CubicSpline}
  \left(\{\bm{u}^{\mathrm{nom}}_j+\bm{\epsilon}_{n,j}\}_{j\in\mathcal{K}}
  \right)(k).
\end{equation}
The actual candidate is clipped componentwise:
\begin{equation}
  \bm{u}^{(n)}_k=clip
  (\tilde{\bm{u}}^{(n)}_k,\bm{u}_{\min},\bm{u}_{\max}).
\end{equation}

The base class also supports independent normal noise at every time step, but
all inspected locomotion configurations select cubic sampling.

\subsection{Cost Normalization and Exponential Weights}

Let $S_n$ be the total rollout cost.  The implementation calculates
\begin{equation}
  S_{\min}=\min_n S_n,\qquad S_{\max}=\max_n S_n,
\end{equation}
then assigns the unnormalized weight
\begin{equation}
  \tilde w_n=\exp\left[
  -\frac{1}{\lambda}
  \frac{S_n-S_{\min}}{S_{\max}-S_{\min}}
  \right].
  \label{eq:weights}
\end{equation}
The new plan is the weighted average of the \emph{absolute candidates}:
\begin{equation}
  \boxed{
  \bm{u}^{+}_k=clip\left(
  \frac{\sum_{n=1}^{N}\tilde w_n\bm{u}^{(n)}_k}
       {\sum_{n=1}^{N}\tilde w_n+10^{-10}},
  \bm{u}_{\min},\bm{u}_{\max}\right).}
  \label{eq:update}
\end{equation}
The default temperature is $\lambda=0.1$ in simulation and $0.15$ in the
walking hardware configuration.

Equation~\eqref{eq:update} is important for interpreting the implementation.
Canonical MPPI is commonly expressed as a weighted perturbation update
$\bm{U}^{+}=\bm{U}+\sum_n w_n\bm{E}^{(n)}$, together with a cost/noise
correction derived from stochastic optimal control.  Here the implementation
instead uses a min--max normalized ranking and directly averages bounded
sampled actions.  It is best described as an MPPI-style sampling controller.

%=============================================================================
\section{Cost Function}
%=============================================================================

\subsection{Quaternion Error}

For predicted and reference quaternions, the scalar sign-invariant error is
\begin{equation}
  d_q(\bm{q},\bm{q}^{\mathrm{ref}})
  =1-\left|\bm{q}^\top\bm{q}^{\mathrm{ref}}\right|.
  \label{eq:quat_error}
\end{equation}
The implementation writes this same scalar into all four quaternion-error
slots.  Hence the effective orientation cost is
\begin{equation}
  \left(\sum_{j=3}^{6}Q_{jj}\right)d_q^2.
\end{equation}
In the default walking configuration, only $Q_{33}=400000$ is nonzero in
that block, so the orientation contribution is $400000d_q^2$.

\subsection{Position and State Tracking}

The original base-position error is removed from the quadratic state error.
Position is instead penalized with a weighted $L_1$ term:
\begin{equation}
  \ell_p(\bm{x})=
  Q_{00}|p_x-p_x^{\mathrm{ref}}|
  +Q_{11}|p_y-p_y^{\mathrm{ref}}|
  +Q_{22}|p_z-p_z^{\mathrm{ref}}|.
  \label{eq:position_cost}
\end{equation}
Let $\bm{e}_x$ equal $\bm{x}-\bm{x}^{\mathrm{ref}}$, with entries
$0{:}3$ set to zero and entries $3{:}7$ replaced by the scalar error in
Eq.~\eqref{eq:quat_error}.  The remaining state cost is
\begin{equation}
  \ell_x=\bm{e}_x^\top\bm{Q}\bm{e}_x+\ell_p.
  \label{eq:state_cost}
\end{equation}
Thus $\ell_x$ includes orientation, gait joint-position tracking, body
velocity tracking, angular-velocity regulation, and gait joint-velocity
tracking.

For the default walking task, the notable diagonal weights are
\begin{equation}
\begin{split}
  \bm{Q}_p &= \diag(300,300,500),\\
  Q_{q,\mathrm{ori}} &=400000,\\
  \bm{Q}_{q_j}&=8000\bm{I}_{12},\\
  \bm{Q}_{v_b}&=\diag(1,10^{-4},10^{-4}),\\
  \bm{Q}_{\omega_b}&=10^{-4}\bm{I}_3,\\
  \bm{Q}_{\dot q_j}&=0.1\bm{I}_{12}.
\end{split}
\end{equation}
These values show that the gait posture and body orientation dominate the
quadratic reference tracking, while forward-speed error receives a much
smaller weight.

\subsection{Virtual PD-Effort Cost}

The cost does not penalize desired position $\bm{u}$ directly.  It constructs
a virtual PD effort
\begin{equation}
  \bm{e}_u=50(\bm{u}-\bm{q}_j)-3\dot{\bm{q}}_j
  \label{eq:virtual_effort}
\end{equation}
and applies
\begin{equation}
  \ell_u=\bm{e}_u^\top\bm{R}\bm{e}_u,
  \qquad \bm{R}=10^{-3}\bm{I}_{12}
\end{equation}
in the default walking configuration.  This is a proxy for actuator effort.
It differs slightly from the rollout actuator in Eq.~\eqref{eq:mjcf_pd},
which uses proportional gain 55 and force saturation.

\subsection{Stage and Rollout Cost}

The stage cost implemented by \code{quadruped\_cost\_np} is
\begin{equation}
  \boxed{
  \ell(\bm{x},\bm{u},\bm{x}^{\mathrm{ref}})
  =\bm{e}_x^\top\bm{Q}\bm{e}_x
  +\bm{e}_u^\top\bm{R}\bm{e}_u
  +\ell_p.}
  \label{eq:stage_cost}
\end{equation}
The rollout cost is the un-discounted sum
\begin{equation}
  S_n=\sum_{k=0}^{H-1}
  \ell(\bm{x}^{(n)}_{k+1},\bm{u}^{(n)}_k,
       \bm{x}^{\mathrm{ref}}_k).
  \label{eq:rollout_cost}
\end{equation}
No explicit slip, contact-force, friction-cone, foot-clearance, self-collision,
control-rate, or terminal term appears in this locomotion cost.  Their effects
are represented only indirectly through dynamics and state/gait tracking.

%=============================================================================
\section{Waypoint and Task Progression}
%=============================================================================

Tasks such as \code{walk\_straight}, \code{walk\_octagon}, \code{stairs}, and
\code{big\_box} define ordered waypoint sequences.  The active waypoint is
considered geometrically reached when
\begin{equation}
  \norm{\bm{p}_g-\bm{p}_b}<r_g,
\end{equation}
where the inspected tasks normally use $r_g=0.2$ m.  The simulator and
hardware loop call \code{next\_goal()} only while this condition holds.

The timer is measured in controller calls rather than seconds.  At the
nominal 100 Hz rate, 100 timer increments correspond to approximately one
second.  When the dwell count completes, \code{next\_goal()} updates the goal
position, planar velocity command, and active gait.  At the last goal it sets
the task-success flag.

The timer does not enforce continuous residence inside the goal region: if
the robot leaves the threshold, accumulated elapsed counts are retained until
the next successful invocation.  This detail matters for interpreting
\code{waiting\_times} as accumulated in-threshold calls rather than a strict
continuous dwell interval.

%=============================================================================
\section{Complete Receding-Horizon Algorithm}
%=============================================================================

\begin{algorithm}[ht]
\caption{\code{legged\_mppi} locomotion update}
\label{alg:legged_mppi}
\begin{algorithmic}[1]
  \REQUIRE Current state $\hat{\bm{x}}_t$, previous plan $\bm{U}_t$,
           active gait and waypoint
  \STATE Compute the goal direction and update the desired body quaternion.
  \STATE Read the next $H$ cyclic gait columns to form
         $\bm{X}^{\mathrm{ref}}_t$.
  \STATE Add Gaussian noise at four knots, interpolate cubic splines, and
         clip $N$ candidate action sequences.
  \STATE Repeat $\hat{\bm{x}}_t$ as the initial state of all candidates.
  \STATE Roll out all $N$ sequences for $H$ MuJoCo steps in parallel.
  \STATE Rotate predicted base linear velocities into their local body frames.
  \STATE Evaluate $S_n$ from position, state, gait, and virtual-effort costs.
  \STATE Min--max normalize $S_n$ and compute exponential weights.
  \STATE Compute $\bm{U}_t^+$ as the weighted mean of absolute candidates.
  \STATE Shift $\bm{U}_t^+$ by one step and hold its tail for warm start.
  \STATE Advance the gait phase by one column.
  \RETURN First desired joint-position command $\bm{u}_{0,t}^+$
\end{algorithmic}
\end{algorithm}

In simulation the first action is applied for one 10 ms MuJoCo step before the
process repeats.  In hardware, the first action is published as twelve
Unitree motor position commands at a nominal 100 Hz ROS loop rate.

%=============================================================================
\section{Default Configuration Summary}
%=============================================================================

\begin{center}
\begin{tabular}{ll}
\toprule
Parameter & Default walking value\\
\midrule
Robot & Unitree Go1, floating base\\
State / input dimension & $37$ / $12$\\
Input interpretation & desired joint position\\
Time step / horizon & $0.01$ s / 40 steps ($0.4$ s)\\
Samples / workers & 30 / 5\\
Temperature & $0.1$ simulation; $0.15$ walking hardware\\
Sampler & cubic spline, four noisy knots\\
Random seed & 42\\
Nominal gait lengths & 74 or 100 samples\\
Contact parameters & \code{o\_solref}$=[0.02,1.0]$\\
Action limits & Go1 joint-position limits\\
Initial plan & fixed crouched joint posture\\
\bottomrule
\end{tabular}
\end{center}

The YAML files for stairs, large-box traversal, and hardware reuse the same
basic solver dimensions while changing cost weights, temperature, noise, or
the environment model.  Task dictionaries additionally select waypoint,
velocity, gait, and model sequences.

%=============================================================================
\section{Comparison with OCS2 Legged MPC}
%=============================================================================

\begin{center}
\begin{tabular}{p{0.22\textwidth}p{0.33\textwidth}p{0.35\textwidth}}
\toprule
Aspect & OCS2 legged MPC & \code{legged\_mppi} locomotion\\
\midrule
Dynamics & Centroidal/SRBD or full centroidal & Full MuJoCo multibody rollout\\
State & Centroidal momentum, base pose, joints & MuJoCo floating-base position and velocity state\\
Optimized input & Contact forces and joint velocities & Desired joint positions\\
Contact timing & Prescribed mode schedule & Joint-space gait prior plus simulated collision\\
Contact force & Explicit decision variable & Internal MuJoCo result\\
Constraints & Mode-dependent equalities, inequalities, barriers & Action clipping, model limits, contact dynamics\\
Solver & SQP, DDP, or IPM & Sampling and exponential weighted averaging\\
Derivatives & Analytical/automatic derivatives & Derivative-free rollout scoring\\
Low-level output & Whole-body controller maps MPC output to torque & Position command sent to simulated or hardware PD servo\\
\bottomrule
\end{tabular}
\end{center}

The two controllers therefore place complexity in different layers.  OCS2
uses an explicit reduced-order OCP and constraint set, followed by a
whole-body controller.  \code{legged\_mppi} places the full robot, actuator,
terrain, and contact response inside the rollout and searches directly in
joint-command space.

%=============================================================================
\section{Implementation Caveats and Limitations}
%=============================================================================

\begin{enumerate}
  \item \textbf{The update differs from canonical MPPI.}
  Cost is normalized by the current batch range, and absolute candidate
  actions are averaged.  The temperature therefore controls selection over
  relative batch cost rather than cost in its original physical units.

  \item \textbf{Equal rollout costs are numerically unsafe.}
  Equation~\eqref{eq:weights} divides by $S_{\max}-S_{\min}$ without an
  epsilon.  If every candidate has the same cost, the weights can become
  non-finite even though the final weight sum contains $10^{-10}$.

  \item \textbf{There is no explicit nominal candidate.}
  Every candidate receives random knot noise.  Consequently, the current
  warm-start plan is not guaranteed to be evaluated unchanged.

  \item \textbf{Gait timing is not a contact constraint.}
  The desired stance/swing behavior can be lost when the joint-space reference
  is not tracked.  Conversely, very large joint weights can make the
  controller behave mainly as a gait playback system.

  \item \textbf{The cost omits several locomotion safety quantities.}
  There is no direct penalty or constraint for foot slip, ground-reaction
  force, friction margin, support polygon, body collision, or joint torque.

  \item \textbf{Virtual and applied effort differ.}
  The cost uses $K_p=50$ and no saturation, while the default MuJoCo actuator
  uses $K_p=55$ and force limits.  Hardware dynamics, latency, and saturation
  add further mismatch.

  \item \textbf{Prediction and execution models can differ.}
  Separate MJCF paths are loaded for the controller and simulator.  Contact
  geometry and solver settings should be audited together for each task.

  \item \textbf{Velocity-frame consistency requires validation.}
  The rollout cost explicitly converts world linear velocity into a body
  frame.  The hardware callback also transforms odometry velocity, but its use
  of the forward rotation should be checked against the frame convention of
  each odometry source.

  \item \textbf{Real-time performance is not guaranteed by the nominal rate.}
  A 100 Hz loop provides only 10 ms per update.  The wall-clock time for 1200
  contact steps, allocation, interpolation, and cost calculation must be
  measured on the target CPU.
\end{enumerate}

%=============================================================================
\section{Recommended Validation}
%=============================================================================

The following measurements connect the formulation to actual closed-loop
behavior:
\begin{enumerate}
  \item mean, 95th, and 99th percentile MPPI update time and 10 ms deadline
        miss rate,
  \item distribution of $S_{\max}-S_{\min}$, effective sample size, and the
        fraction of non-finite weights,
  \item base position, height, quaternion error, and body-frame velocity error,
  \item joint/gait tracking error, actuator saturation rate, and measured
        effort,
  \item foot contact sequence, ground-reaction force, slip velocity, and body
        collision events,
  \item one-step prediction error between the rollout model and simulator or
        hardware observation, and
  \item waypoint residence time, transition count, and task-completion rate.
\end{enumerate}

Useful ablations include normal versus cubic noise, the current absolute-action
average versus a perturbation update, inclusion of a noise-free candidate,
removal of min--max normalization, varying the gait-state weights, and using
identical versus deliberately mismatched prediction/execution contact models.

%=============================================================================
\section{Pipeline Summary}
%=============================================================================

\begin{equation*}
\boxed{
\begin{array}{c}
\text{Current Go1 state }\hat{\bm{x}}_t\\
\downarrow\\
\text{Waypoint body target + phase-aligned joint gait}\\
\downarrow\\
\text{Warm-start desired joint-position sequence}\\
\downarrow\\
\text{Cubic Gaussian sampling }(30\times40\times12)\\
\downarrow\\
\text{Parallel full-dynamics MuJoCo rollouts}\\
\downarrow\\
\text{Body/gait tracking + virtual PD-effort cost}\\
\downarrow\\
\text{Min--max normalization + exponential weighting}\\
\downarrow\\
\text{Execute the first desired joint-position action}\\
\downarrow\\
\text{Shift the plan and advance the gait phase}
\end{array}}
\end{equation*}

The defining design choice is to optimize desired joint-position trajectories
through full MuJoCo contact dynamics.  A periodic gait supplies a strong
motion prior, while sampling searches locally for whole-body actions that
reduce pose, velocity, gait, and virtual-effort cost in receding-horizon
feedback.

\end{document}
